
\documentclass{beamer}
\usepackage[utf8]{inputenc}
\usepackage[T1]{fontenc}
\usepackage[french]{babel}
\usetheme{Madrid}
\usepackage{xcolor}
% Couleurs personnalisées
\definecolor{myblue}{HTML}{1E90FF}
\definecolor{mygreen}{HTML}{27ae60}
\definecolor{myorange}{HTML}{e67e22}
\definecolor{mypurple}{HTML}{8e44ad}
\definecolor{myred}{HTML}{c0392b}
\setbeamercolor{frametitle}{bg=myblue!10,fg=myblue!80!black}
\setbeamercolor{block title}{bg=mygreen!20,fg=mygreen!80!black}
\setbeamercolor{block body}{bg=mygreen!5,fg=black}
\setbeamercolor{item}{fg=myblue!70!black}
\setbeamertemplate{itemize item}{\color{myblue}$\blacktriangleright$}
\setbeamertemplate{itemize subitem}{\color{myorange}$\star$}

\title{Gestion de projet informatique :\newline méthodes traditionnelles vs agiles}
\author{Anraud Gadille, Maël Mignard}
\date{\today}

\begin{document}

\begin{frame}
  \titlepage
\end{frame}

\section{Introduction : Projet informatique}
\begin{frame}{\color{myblue}Définition d'un projet informatique}
  \scriptsize
  \begin{columns}
    \column{0.48\textwidth}
    {\setbeamercolor{block title}{bg=myblue!20,fg=myblue!80!black}
    \setbeamercolor{block body}{bg=myblue!5,fg=black}
    \begin{block}{Concepts}
      \begin{itemize}
        \item \textbf{Objectif} : Créer ou améliorer un logiciel/système.
        \item \textbf{Phases} : Analyse $\rightarrow$ Conception $\rightarrow$ Développement $\rightarrow$ Tests $\rightarrow$ Déploiement $\rightarrow$ Maintenance.
        \item \textbf{Ressources} : Temps, budget, équipe, outils.
        \item \textbf{Risques} : Retards, dépassement de budget, changements.
        \item \textbf{Livrable} : Résultat final fonctionnel.
      \end{itemize}
    \end{block}}
    \column{0.48\textwidth}
    {\setbeamercolor{block title}{bg=myorange!20,fg=myorange!80!black}
    \setbeamercolor{block body}{bg=myorange!5,fg=black}
    \begin{block}{Exemples}
      \begin{itemize}
        \item Site web pour vendre des gâteaux.
        \item Discuter $\rightarrow$ Dessiner $\rightarrow$ Coder $\rightarrow$ Tester $\rightarrow$ Lancer.
        \item 3 mois, 2000€, développeur + designer.
        \item Développeur malade $\rightarrow$ retard.
        \item Site en ligne avec panier et paiement sécurisé.
      \end{itemize}
    \end{block}}
  \end{columns}
\end{frame}

\begin{frame}{\color{mygreen}Bases, acteurs et outils}
  \scriptsize
  \begin{columns}
    \column{0.48\textwidth}
    {\setbeamercolor{block title}{bg=mygreen!20,fg=mygreen!80!black}
    \setbeamercolor{block body}{bg=mygreen!5,fg=black}
    \begin{block}{Bases}
      \begin{itemize}
        \item Planifier, exécuter, suivre et clôturer un projet pour atteindre un objectif précis.
      \end{itemize}
    \end{block}}
    {\setbeamercolor{block title}{bg=mypurple!20,fg=mypurple!80!black}
    \setbeamercolor{block body}{bg=mypurple!5,fg=black}
    \begin{block}{Acteurs}
      \begin{itemize}
        \item Chef de projet
        \item Équipe
        \item Clients
      \end{itemize}
    \end{block}}
    \column{0.48\textwidth}
    {\setbeamercolor{block title}{bg=myorange!20,fg=myorange!80!black}
    \setbeamercolor{block body}{bg=myorange!5,fg=black}
    \begin{block}{Outils}
      \begin{itemize}
        \item Logiciels pour organiser, collaborer et suivre l’avancement
        \item Planification
        \item Communication
        \item Suivi
      \end{itemize}
    \end{block}}
  \end{columns}
\end{frame}

\begin{frame}{\color{myorange}Méthodes traditionnelles}
  \scriptsize
  \begin{columns}
    \column{0.48\textwidth}
    {\setbeamercolor{block title}{bg=myorange!20,fg=myorange!80!black}
    \setbeamercolor{block body}{bg=myorange!5,fg=black}
    \begin{block}{Phases}
      \begin{itemize}
        \item Initiation
        \item Planification
        \item Exécution
        \item Suivi et Contrôle
        \item Clôture
      \end{itemize}
    \end{block}}
    \column{0.48\textwidth}
    {\setbeamercolor{block title}{bg=myred!20,fg=myred!80!black}
    \setbeamercolor{block body}{bg=myred!5,fg=black}
    \begin{block}{Objectifs}
      \begin{itemize}
        \item Définir les objectifs, le périmètre, les parties prenantes et la faisabilité du projet.
        \item Établir un plan détaillé : tâches, calendrier, budget, ressources et risques.
        \item Réaliser les tâches prévues selon le plan.
        \item Surveiller l’avancement, comparer avec le plan et corriger les écarts si nécessaire.
        \item Finaliser le projet, livrer le produit et effectuer un bilan (retour d’expérience).
      \end{itemize}
    \end{block}}
  \end{columns}
\end{frame}
\begin{frame}{\color{mypurple}Autres méthodes traditionnelles : Cycle en V vs PERT}
  \scriptsize
  \begin{columns}
    \column{0.48\textwidth}
    {\setbeamercolor{block title}{bg=mypurple!20,fg=mypurple!80!black}
    \setbeamercolor{block body}{bg=mypurple!5,fg=black}
    \begin{block}{Cycle en V}
      \begin{itemize}
        \item \textbf{Définition} : Modèle séquentiel, chaque étape doit être validée avant la suivante.
        \item \textbf{Approche} : Séquentielle, structurée et peu flexible.
        \item \textbf{Planification} : Précise et figée dès le départ.
        \item \textbf{Gestion des risques} : Identification anticipée, peu d'adaptations en cours de projet.
        \item \textbf{Utilisation} : Projets avec exigences stables et bien définies.
        \item \textbf{Avantages} : Étapes claires, suivi facile, documentation complète.
        \item \textbf{Inconvénients} : Manque de souplesse, difficile à adapter aux changements.
      \end{itemize}
    \end{block}}
    \column{0.48\textwidth}
    {\setbeamercolor{block title}{bg=myblue!20,fg=myblue!80!black}
    \setbeamercolor{block body}{bg=myblue!5,fg=black}
    \begin{block}{PERT}
      \begin{itemize}
        \item \textbf{Définition} : Méthode de gestion axée sur l’ordonnancement et la planification des tâches.
        \item \textbf{Approche} : Réseau de tâches, gestion des dépendances et des incertitudes.
        \item \textbf{Planification} : Dynamique, estimation des durées et identification des chemins critiques.
        \item \textbf{Gestion des risques} : Évolutive, ajustements possibles en fonction de l’avancement.
        \item \textbf{Utilisation} : Projets complexes avec de nombreuses tâches et dépendances.
        \item \textbf{Avantages} : Optimisation des délais, meilleure gestion des imprévus.
        \item \textbf{Inconvénients} : Mise en œuvre complexe, nécessite des outils adaptés.
      \end{itemize}
    \end{block}}
  \end{columns}
\end{frame}

\section{Méthodes agiles}
\begin{frame}{Méthodes agiles}
  \scriptsize
  \textbf{Scrum}
  \begin{itemize}
    \item \textbf{Fonctionnement} : Étapes courtes (sprints) + réunions quotidiennes.
    \item \textbf{Qui participe ?} : Product Owner, Scrum Master, équipe
    \item \textbf{Durée} : 2-4 semaines/étape
    \item \textbf{Avantages} : Flexible, livraisons fréquentes.
    \item \textbf{Inconvénients} : Équipe expérimentée nécessaire.
    \item \textbf{Pour quel projet ?} : Logiciels, nouveaux produits.
  \end{itemize}
  \vspace{0.3cm}
  \textbf{Kanban}
  \begin{itemize}
    \item \textbf{Fonctionnement} : Tableau visuel (À faire/En cours/Terminé).
    \item \textbf{Qui participe ?} : Équipe auto-organisée
    \item \textbf{Durée} : Continu
    \item \textbf{Avantages} : Simple, adaptable.
    \item \textbf{Inconvénients} : Moins structuré pour gros projets.
    \item \textbf{Pour quel projet ?} : Maintenance, tâches répétitives.
  \end{itemize}
  \vspace{0.3cm}
  \textbf{XP (Extreme Programming)}
  \begin{itemize}
    \item \textbf{Fonctionnement} : Qualité du code + tests automatiques.
    \item \textbf{Qui participe ?} : Développeurs + client
    \item \textbf{Durée} : 1-2 semaines/étape
    \item \textbf{Avantages} : Code solide, peu de bugs.
    \item \textbf{Inconvénients} : Équipe technique requise.
    \item \textbf{Pour quel projet ?} : Projets complexes, qualité critique.
  \end{itemize}
\end{frame}

\begin{frame}{Comparatif : Méthode Agile vs Traditionnelle}
  \scriptsize
  \begin{columns}
    \column{0.48\textwidth}
    \textbf{Méthode Agile}
    \begin{itemize}
      \item \textbf{Flexibilité} : Très flexible, adaptations fréquentes possibles.
      \item \textbf{Livraison \& implication client} : Livraisons fréquentes (2-4 semaines), feedback client régulier.
      \item \textbf{Documentation} : Légère, axée sur le produit fonctionnel.
      \item \textbf{Risques} : Réduction des risques grâce aux ajustements réguliers et au feedback client.
    \end{itemize}
    \column{0.48\textwidth}
    \textbf{Méthode Traditionnelle}
    \begin{itemize}
      \item \textbf{Flexibilité} : Peu flexible, changements difficiles après le démarrage.
      \item \textbf{Livraison \& implication client} : Livraison unique à la fin, validation client surtout au début et à la fin.
      \item \textbf{Documentation} : Lourde, chaque phase est documentée en détail.
      \item \textbf{Risques} : Risques élevés si les besoins initiaux sont mal définis (découverte tardive des problèmes).
    \end{itemize}
  \end{columns}
\end{frame}



\end{document}
